\section{Discussion and limitations}
The experiments support a distinction between identifying evidence, authorizing
its spatial effect, and recovering a complete task output. Identity admission
and complete-membership controls show that these capabilities affect accuracy;
reconstruction and response-decision controls show that better source agreement
alone is insufficient. Their interventions and cohorts differ, so their gains
cannot be added. The contribution is the evidence-to-response design and its
measured behavior, built with native correspondence and established fusion
operations.

The common contract has typed implementations. Current evidence does not show
one admission rule or selector improving every task and teacher configuration.
The L11 null constraint has no demonstrated accuracy effect. The benchmark
cohort has informed development, and new NVOS variants are development
experiments rather than independent final tests. More teacher configurations
and unseen scenes are needed for broad transfer claims.

Source access, teacher strength and carrier geometry differ from some published
baselines. Gaussian versus mesh score differences are not method gains. LERF
selection uses projected GT, UQIS is a local protocol, and original NVOS v10 differs from
LUDVIG in target-RGB access and GT boundaries. The image-conditioned extension
explicitly changes that access, but keeps source-trained geometry and produces
view-conditioned updates. Early static image refinement improved IoU while
reducing AP; completing the original signed reference response subsequently
improves both. The N12--N14 variants nevertheless increase residuals against
their original source teachers on all eight tasks. Target-view improvements therefore
do not establish a universally improved or query-global 3D field.

The controls localize three distinct limitations. First, memory transport of
an incomplete original reference mask does not restore missing query scope.
Completing that reference before transport recovers the missing orchid flowers.
Second, source annotation agreement does not establish either complete extent
or accurate target boundaries: unconditional extra prompting can enlarge an
already adequate reference response. Source feasibility and predicted quality
are thus measured separately, with fixed decisions for every task. Third,
native-ray fitting reduces transport loss but cannot remove false positives
already present in the image response; fern and trex still expose internal
background leakage and thin-structure loss. N15 tests cross-view consistency:
completed-reference source-only fusion recovers 92.7461 IoU, but equally
weighting these sources with target evidence lowers N14's IoU by 0.1428 points
while increasing AP by 0.0346. Completed-source Huber residual falls from
0.004626 to 0.000748, yet foreground recall also falls. This controlled result
separates source-teacher consistency from task accuracy. The fusion branch is
stopped without a post-hoc weight search; it does not establish a universal
performance ceiling. The unchanged v10 threshold oracle (90.1098) diagnoses only that
original score field, not an upper bound on new spatial responses.

N16 shows that one extra target decode is also insufficient as a uniform
repair rule. Its local control modestly improves thresholded IoU but lowers AP;
the tradeoff already appears in direct image responses, before 3D fitting.
Its mean-IoU gain over N14 disappears when the correlated horns queries are
omitted together. This descriptive sensitivity is not a new validation split.
The new maximum IoU is retained in the tables rather than hidden, but the
extra decoding stage is not required by the recommended method. These stopped
branches support an empirical stopping decision for the current refinement
recipes, not a claim that no stronger evidence or architecture can help.
The subsequent A24 numerical audit exactly replays all 24 image responses.
Evaluating the same probability interpolation with float64 and a stable
complement-based ranking changes mean AP by less than 0.002 points and changes
no hard masks. The local image response remains 2.2124 AP points below matched
N14. Thus the suspected final sigmoid/interpolation precision effect does not
explain the observed ranking loss. This audit does not recompute 3D results.

The N11 geometric and v10-response anchors were nearly tied, so a distinct
geometric-admission advantage was not established. N13's paired probability
residual loses IoU on all eight tasks against its matched complete replacement.
Its exact no-intervention property is therefore insufficient evidence for a
better response rule. These alternatives remain diagnostic controls, not
mandatory stages of the current reference-completion pipeline. New admission
rules were developed after inspecting N12; they are explicitly development
results, without scene-name routing or target-score selection at inference.
Independent scorer checks reduce
known implementation concerns but do not prove that every possible bias has
been excluded.

\paragraph{Reproducibility.}
Versioned recipes, source/input hashes, per-task scores, protocol audits and
component controls bind claims to code and outputs. Teachers, maps and geometry
are referenced in place. New dense predictions are bounded and transient;
source snapshots, small models and reports persist. Table generation does not
run inference. Incomplete historical source archives remain explicitly marked.

\section{Conclusion}
Query-conditioned evidence fusion connects native evidence to geometric
responses while preserving distinct semantic and complete-selection outputs.
Our v10 system supplies the five-task implementation and main measurements.
Matched controls identify effective identity, response-use and extent operations;
negative controls show that better source fit alone is insufficient. NVOS
evidence experiments and optional response replacement are specific extensions,
with their effects separated from established method capability.
